% ============================================================
% Methodology excerpt — AI-assisted functional test suite for
% a legacy Fortran earth system model
% ============================================================

\subsection{AI-Assisted Functional Test Suite Development}

Developing a functional test suite for a large, legacy Fortran code base presents
challenges that are representative of the broader class of Earth System Models
(ESMs): subroutines are often tightly coupled to model-specific derived types,
access modifiers frequently restrict calling from external programs, and the sheer
number of subroutines makes exhaustive manual coverage prohibitive.  We describe
here a methodology in which an AI coding assistant (Claude, Anthropic) was used
interactively to design, implement, and expand a functional test suite for
CLM-ml~v2, a multilayer plant-canopy extension of the Community Land Model.  The
methodology is organized around three interlocking concerns: \emph{testability
assessment}, \emph{driver architecture}, and \emph{verification strategy}, each
of which generalizes to any Fortran ESM that uses similar compilation and linking
conventions.

\subsubsection{Testability Assessment and Subroutine Classification}

The first step in the methodology is systematic classification of all public and
private subroutines according to their testability.  The AI was tasked with reading
each Fortran source module and categorizing every subroutine into one of three
tiers.  \emph{Tier~1} subroutines take purely scalar or fixed-size array arguments
(real, integer, or logical) with no dependency on model-specific derived types, and
can therefore be called directly from a standalone test program with no further
modification to the source.  Examples from this work include mathematical utilities
(\texttt{quadratic}, \texttt{tridiag}, gamma and beta functions), thermodynamic
functions (\texttt{SatVap}, \texttt{LatVap}), and photosynthesis temperature-response
factors (\texttt{ft}, \texttt{fth}, \texttt{fth25}).  \emph{Tier~2} subroutines
contain self-contained, mathematically isolable logic but are embedded in loops over
a large model-specific derived type (in this case \texttt{mlcanopy\_type}) or are
declared \texttt{private} to their module.  These require either a visibility change
(\texttt{public ::} declaration) or the addition of a \emph{scalar helper}
subroutine—a new, minimal public routine that extracts the pure computation and
accepts only primitive scalar arguments.  \emph{Tier~3} subroutines are deeply
coupled to the full model state (multiple derived types, file I/O, or MPI
communication) and cannot be meaningfully tested without standing up a complete model
column; these were noted as out of scope for unit-level functional testing.  This
three-tier classification can be performed by any AI assistant capable of reading
Fortran source files and should be the first step for any ESM test-suite project,
because it bounds the effort and prevents wasted work on subroutines whose coupling
cost exceeds the testing benefit.

\subsubsection{Source Modifications to Expose Testable Interfaces}

Where Tier~2 subroutines were identified as scientifically important, two types of
minimal, non-breaking source modifications were made.  The first type is a simple
visibility change: adding a \texttt{public~::} declaration for a subroutine that
already has a scalar-friendly interface but was declared \texttt{private} for
encapsulation reasons.  In this work, seven subroutines in
\texttt{MLLeafPhotosynthesisMod} and \texttt{MLCanopyTurbulenceMod} were exposed
this way, including the Monin--Obukhov stability functions
(\texttt{phim\_monin\_obukhov}, \texttt{psim\_monin\_obukhov}, and their scalar
counterparts) and the canopy-turbulence parameters \texttt{GetBeta},
\texttt{GetPrSc}, and \texttt{GetPsiRSL}.  The second type is the \emph{scalar
helper pattern}: a new public subroutine is added to the module, immediately
following the original, that accepts the same mathematical inputs as scalar
arguments, performs the identical computation, and returns scalar outputs.  The
original subroutine is left entirely unchanged; the helper merely lifts the
relevant formulas out of the loop-over-derived-type structure.  Four such helpers
were introduced: \texttt{CalcWettedFraction} (for canopy interception physics),
\texttt{CalcLeafHeatCapacity} (for leaf energy storage), \texttt{CalcLeafBoundaryLayer}
(for boundary-layer conductances), and \texttt{CalcNitrogenScale} (for the
exponential canopy nitrogen profile).  Because all modifications are purely additive
and introduce no new execution paths in existing subroutines, they carry negligible
risk of altering model behavior and can be merged into any production code base
without review overhead.  This pattern generalizes directly to any Fortran ESM: the
key criterion is that the subroutine's physics depends only on its arguments, not on
global state that cannot be constructed outside the full model context.

\subsubsection{Two-Layer Driver Architecture}

The test framework rests on a two-layer architecture that separates the
Fortran execution environment from the Python testing logic.  In the first layer,
each subroutine under test receives a dedicated \emph{Fortran driver program}—a
minimal \texttt{program} unit that reads a Fortran namelist from standard input,
calls the target subroutine with those inputs, and writes the outputs to standard
output as \texttt{KEY=VALUE} pairs in IEEE~754 double-precision exponential format
(\texttt{ES24.16}).  Namelists were chosen because they are a native Fortran
mechanism requiring no external libraries, support mixed scalar and array inputs
without custom parsing, and are trivial to generate programmatically from Python.
The \texttt{KEY=VALUE} output convention requires no third-party output libraries
and is unambiguously parseable by a single regular expression.  Each driver is
compiled and linked against the full set of object files produced by the main model
build (excluding only the model's own \texttt{main} program), so the test executable
uses exactly the same compiled physics as the production model.  A single GNU
\texttt{make} pattern rule handles all 26 drivers with no per-driver configuration.
In the second layer, a Python \texttt{pytest} suite provides the testing logic.  A
shared utility function, \texttt{run\_fortran(exe\_name, inputs\_dict)}, constructs
the namelist string, pipes it to the driver executable via \texttt{subprocess.run},
and returns a dictionary mapping output key names to floating-point values.  This
clean interface means that every Python test is free of Fortran syntax and can
express assertions in idiomatic Python with standard \texttt{assert} statements.  A
session-scoped \texttt{pytest} fixture in \texttt{conftest.py} invokes \texttt{make
all} once at the start of each test session, so the executables are always
up-to-date before any test runs.  This architecture is directly portable to any
Fortran ESM that uses a batch-oriented build system (Make, CMake, or similar),
because the only requirement is that object files from the main build can be linked
with an external \texttt{program} unit.

\subsubsection{Verification Strategy: Property Tests and Golden-File Regression}

Each Python test module combines two complementary assertion strategies.
\emph{Property-based tests} verify mathematical or physical invariants that hold
for all valid inputs and can be derived analytically from the governing equations.
Examples include: the Clausius--Clapeyron relation (saturation vapor pressure is
monotonically increasing with temperature and its derivative is positive); Vieta's
formulas for the roots of a quadratic; the symmetry $B(a,b)=B(b,a)$ and boundary
conditions $F(0;a,b)=0$, $F(1;a,b)=1$ of the beta distribution; the neutral
condition $\phi_m(0)=1$ and stable asymptote $\phi_m(\zeta)=1+5\zeta$ of the
Monin--Obukhov stability functions; the implicit consistency equation
$\beta\,\phi_m(L_{cL}\beta^2) = \beta_\mathrm{neutral}$ for the canopy wind
attenuation ratio; the exact cancellation $\psi_m = v_{kc}/\beta$, $\psi_c = 0$
at $z_a = h_c$ in the Roughness Sublayer formulation; and the Butcher tableau
consistency condition $\sum_j b_j = 1$ for the Runge--Kutta integration scheme.
\emph{Reference-value tests} compare computed outputs against tabulated or
analytically computed expected values at specific inputs, providing a precise
numerical anchor.  The combination of the two strategies means that a regression
is caught whether it manifests as a qualitative sign error (property failure) or a
subtle numerical drift (reference-value failure).

In addition to the per-subroutine tests, a \emph{golden-file regression system}
captures the complete input--output record of a trusted model build.  A script,
\texttt{generate\_golden.py}, runs every driver executable against a curated set
of representative input cases, recording inputs and outputs as JSON files in a
committed \texttt{tests/golden/} directory.  A dedicated test module,
\texttt{test\_golden.py}, loads these files at every subsequent \texttt{pytest}
session and re-runs all cases, asserting that outputs match the golden record
within a tight numerical tolerance.  This provides a safety net for unintentional
regressions introduced by compiler upgrades, optimization flag changes, or
algorithm edits.  The golden files are committed alongside the source code, so
the regression baseline travels with the code in version control.  When a change
is \emph{intentional}, the developer re-runs \texttt{generate\_golden.py} and
commits the updated JSON files together with the source change, making the audit
trail explicit.  This two-stage approach—\emph{per-subroutine property/reference
tests} for physical correctness plus \emph{golden-file regression} for
reproducibility—provides defense in depth and is directly applicable to any ESM
with a deterministic, bit-reproducible build.

\subsubsection{Parent-Level Integration Testing}

A third test pattern was developed for subroutines that contain \emph{internal
analytical consistency assertions}—cases where the subroutine itself checks whether
its own outputs satisfy a mathematical identity and aborts if they do not.  These
assertions cannot be exercised by a single-layer scalar test alone, because they
involve a sum over all canopy layers.  For \texttt{CanopyNitrogenProfile}, the
relevant identity is
\begin{equation}
  \sum_{i=1}^{N} V_{\mathrm{cmax},25}^{(i)}\,\Delta\mathrm{PAI}^{(i)}
  \;=\;
  V_{\mathrm{cmax},25}^{\mathrm{top}}
  \,\frac{1 - e^{-k_n\,\mathrm{PAI}_{\mathrm{tot}}}}{k_n},
  \label{eq:nitrogen_integral}
\end{equation}
which states that the discrete layer-by-layer sum of the exponential nitrogen
decay profile must reproduce the closed-form analytical integral over the entire
canopy.  A \emph{parent-level integration driver} was written that accepts a full
multi-layer canopy configuration via namelist, replicates the production loop using
the scalar helper \texttt{CalcNitrogenScale}, and outputs both the numerical sum
and the analytical reference, allowing the Python test to verify Eq.\
\eqref{eq:nitrogen_integral} for arbitrary layer counts and PAI distributions.
Critically, the test also verifies \emph{layer-count invariance}: discretizing the
same total canopy depth into 3 layers or 6 layers yields identical integrated
$V_{\mathrm{cmax},25}$, confirming that the integration scheme is free of
discretization error.  This pattern generalizes to any ESM subroutine that
contains an internal energy, mass, or tracer conservation check: the check provides
a natural test oracle, and the parent-level driver simply needs to supply the
structured inputs required to trigger it.

\subsubsection{Iterative AI-Driven Expansion}

The test suite was built through multiple rounds of interaction in which the AI was
asked first to \emph{assess} remaining testability gaps and then to
\emph{implement} the identified tests.  Each assessment round involved the AI
reading the current source modules, comparing them against the list of already-tested
subroutines, classifying the remainder by tier, and providing a ranked list of
candidates with rationale.  This iterative pattern—assess, implement, reassess—is
well-suited to AI assistance because the assessment step requires broad but shallow
reading of many files (a task at which large language models are efficient) while
the implementation step requires precise, consistent application of a fixed pattern
(the driver template, the Python test structure, the Makefile rule) to each new
candidate.  The boundary between what is tractable and what is not (Tier~3
subroutines, root-finding functions that require a model-specific function pointer)
was surfaced by the AI during assessment and accepted without further implementation
effort, avoiding wasted work.  Over six implementation rounds, the suite grew from
15 to 26 executables covering 10 modules and 32 distinct subroutines or functions.
The total number of lines of test code (Fortran drivers, Python tests,
infrastructure) produced during these rounds exceeds 4,000, all of which was
verified for correctness against known analytical results or physical constraints
before being committed.  We expect this iterative, AI-driven workflow to be
reproducible for any Fortran ESM: the key inputs are the module source files and
a description of the build system; the key output is a growing library of
executable, self-describing tests that remain useful long after the original
development effort concludes.
